\paragraph{Задачи}

\subparagraph{2.1. Сортировка вставкой малых массивов в процессе сортировки слиянием}

Несмотря на то что с увеличением количества сортируемых элементов время сортировки методом слияния в наихудшем случае растёт как $\Theta(n \lg n)$,
а время сортировки вставками --- как $\Theta(n^2)$, благодаря постоянным множителям
на практике для малых размеров задач на большинстве машин сортировка вставками выполняется быстрее.
Таким образом, имеет смысл использовать сортировку вставками в процессе сортировки методом слияния, когда подзадачи становятся достаточно маленькими.
Рассмотрим модификацию алгоритма сортировки слиянием,
в котором $n/k$ подмассивов длиной $k$ сортируются вставками, после чего они объединяются с помощью обычного механизма слияния.
Величина $k$ должна быть найдена в процессе решения задачи.

\begin{enumerate}
    \item[а.] Покажите, что сортировка вставками позволяет отсортировать $n/k$ подпоследовательностей длиной $k$ каждая за время $\Theta(nk)$ в худшем случае.

    \textbf{Ответ:}
    В наихудшем случае сортировка вставками сортирует подпоследовательность длиной $k$ за время $\Theta(k^2)$.
    Поскольку таких подпоследовательностей $n/k$, общее время сортировки вставками составит $\Theta(nk)$.

    \item[б.] Покажите, как выполнить слияние этих подпоследовательностей за время $\Theta(n \lg(n/k))$ в наихудшем случае.

    Для одного уровня слияния требуется время $\Theta(n)$.
    Если построить дерево рекурсии, то на каждом уровне общее количество объединяемых элементов равно $n$.
    Количество уровней равно $\Theta(\lg(n/k))$.
    Следовательно, общее время слияния составляет $\Theta(n \lg(n/k))$.

    \item[в.] Если такой модифицированный алгоритм выполняется за время $\Theta(nk + n \lg(n/k))$ в наихудшем случае,
    то чему равно наибольшее значение $k$ как функции от $n$, при котором модифицированный алгоритм в $\Theta$-обозначениях имеет то же время работы,
    что и стандартная сортировка слиянием?

    Для этого приравняем времена работы алгоритмов:
\[
    c_1(nk + n \lg(n/k)) = c_2 n \lg n
\]
\[
    c_1 k + c_1 \lg(n/k) = c_2 \lg n
\]
\[
    c_1 k + c_1 \lg n - c_1 \lg k = c_2 \lg n
\]
\[
    c_1 k - c_1 \lg k = c_2 \lg n - c_1 \lg n
\]
\[
    c_1 (k - \lg k) = (c_2 - c_1)\lg n
\]

    Из равенства
\[
    k - \lg k = \Theta(\lg n)
\]
заметим, что $\lg k = o(k)$, поэтому
\[
    k - \lg k = \Theta(k)
\]
Следовательно,
\[
    k = \Theta(\lg n).
\]

    \item[г.] Как следует выбирать $k$ на практике?

    На практике выбор $k$ можно осуществлять на основе эмпирических данных, полученных при замерах времени работы алгоритмов.
    Можно найти такое значение $n$, при котором сортировка слиянием и вставками показывает одинаковое время работы,
    и использовать значение, немного меньшее найденного, в качестве порогового для переключения на сортировку вставками.
\end{enumerate}

\subparagraph{2.2. Корректность пузырьковой сортировки}

Пузырьковая сортировка представляет собой популярный, но неэффективный алгоритм сортировки. 
В его основе лежит многократная перестановка соседних элементов, нарушающих порядок сортировки.

\begin{verbatim}
BUBBLESORT(A)
1  for i = 1 to A.length - 1
2      for j = A.length downto i + 1
3          if A[j] < A[j - 1]
4              поменять A[j] и A[j - 1] местами
\end{verbatim}

\begin{enumerate}
    \item[а.] Пусть $A'$ обозначает выход процедуры $\text{BUBBLESORT}(A)$.
    Для доказательства корректности процедуры $\text{BUBBLESORT}$ необходимо доказать, что она завершается и что
    \[
    A'[1] \leq A'[2] \leq \dots \leq A'[n],
    \]
    где $n = A.\text{length}$.
    Что ещё необходимо доказать для того, чтобы показать, что процедура $\text{BUBBLESORT}$ действительно выполняет сортировку?

    \textbf{Ответ:} \\
    Для подтверждения корректности алгоритма необходимо также доказать,
    что результирующий массив $A'$ является перестановкой исходного массива $A$.

    В следующих двух частях доказывается неравенство выше.

    \item[б.] Точно сформулируйте инвариант цикла \texttt{for} в строках 2--4 и докажите, что он выполняется.

    \textbf{Инвариант цикла:} \\
    На каждой итерации цикла при текущем значении $j$
    элемент $A[j]$ является наименьшим среди элементов $A[j \ldots n]$.

    \textbf{Инициализация:} \\
    При $j = A.\text{length}$ подмассив $A[j \ldots n]$ содержит один элемент $A[j]$,
    который является наименьшим.

    \textbf{Сохранение:} \\
    Предположим, что перед итерацией при текущем значении $j$
    элемент $A[j]$ является наименьшим среди $A[j \ldots n]$.
    Сравниваются элементы $A[j]$ и $A[j-1]$, и при необходимости выполняется обмен,
    в результате чего на позиции $j-1$ оказывается наименьший элемент из $A[j-1 \ldots n]$.
    После уменьшения $j$ инвариант сохраняется.

    \textbf{Завершение:} \\
    При завершении цикла $j = i$, и по инварианту элемент $A[i]$
    является наименьшим среди элементов $A[i \ldots n]$.

    \item[в.] Сформулируйте инвариант внешнего цикла.

    \textbf{Инвариант цикла:} \\
    На каждой итерации цикла \texttt{for} в строках 1--4 подмассив $A[1 \ldots i-1]$
    содержит наименьшие элементы массива и отсортирован по возрастанию.

    \textbf{Инициализация:} \\
    При $i = 1$ подмассив $A[1 \ldots i-1]$ пуст, следовательно, он тривиально отсортирован.

    \textbf{Сохранение:} \\
    Предположим, что перед итерацией при текущем значении $i$
    подмассив $A[1 \ldots i-1]$ содержит наименьшие элементы массива и отсортирован.
    Внутренний цикл перемещает наименьший элемент из подмассива $A[i \ldots n]$
    на позицию $i$.
    Таким образом, после увеличения $i$ подмассив $A[1 \ldots i]$
    содержит наименьшие элементы и остаётся отсортированным.

    \textbf{Завершение:} \\
    При завершении цикла $i = n$,
    подмассив $A[1 \ldots n]$ содержит все элементы массива в отсортированном порядке.

    \item[г.] Определите время работы пузырьковой сортировки в наихудшем случае.

    \textbf{Ответ:} \\
    Время работы можно оценить следующим образом:
    \[
    \sum_{i=1}^{n-1} \sum_{j=i+1}^{n} 1
    = \sum_{i=1}^{n-1} (n - i)
    = \frac{n(n-1)}{2}
    = \Theta(n^2).
    \]
\end{enumerate}

\subparagraph{2.3. Корректность правила Горнера}

Следующий фрагмент кода реализует правило Горнера для вычисления полинома
\[
\begin{aligned}
P(x) &= \sum_{k=0}^{n} a_k x^k \\
&= a_0 + x(a_1 + x(a_2 + \dots + x(a_{n-1} + x a_n) \dots))
\end{aligned}
\]
для заданных коэффициентов $a_0, a_1, \dots, a_n$ и значения $x$.

\medskip
\noindent
\begin{tabular}{l l}
1 & $y = 0$ \\
2 & \textbf{for} $i = n$ \textbf{downto} 0 \\
3 & \quad $y = a_i + x \cdot y$
\end{tabular}
\medskip

\begin{itemize}
    \item[\textbf{а.}] Чему равно время работы этого фрагмента кода правила Горнера в $\Theta$-обозначениях?
    
    \textbf{Ответ:} \\
    $\Theta(n)$

    \item[\textbf{б.}] Напишите псевдокод, реализующий алгоритм обычного вычисления полинома, когда каждое слагаемое полинома вычисляется отдельно. Определите асимптотическое время работы этого алгоритма и сравните его со временем работы алгоритма, основанного на правиле Горнера.
    
    \item[\textbf{в.}] Рассмотрим следующий инвариант цикла.
    
    В начале каждой итерации цикла \textbf{for} в строках 2 и 3
    \[
    y = \sum_{k=0}^{n-(i+1)} a_{k+i+1} x^k \text{.}
    \]
    Рассматривайте сумму без членов как равную нулю. Следуя структуре доказательства инварианта цикла, которая использовалась ранее в данной главе, воспользуйтесь указанным инвариантом цикла, чтобы показать, что по завершении работы $y = \sum_{k=0}^n a_k x^k$.
    
    \item[\textbf{г.}] Сделайте заключение, что в приведенном фрагменте кода правильно вычисляется значение полинома, который задается коэффициентами $a_0, a_1, \dots, a_n$.
\end{itemize}

\subparagraph{2.4. Инверсии}

Пусть $A[1 \dots n]$ представляет собой массив из $n$ различных чисел. Если $i < j$ и $A[i] > A[j]$, то пара $(i, j)$ называется \textbf{инверсией} $A$.

\begin{itemize}
    \item[\textbf{а.}] Перечислите пять инверсий массива $\langle 2, 3, 8, 6, 1 \rangle$.
    
    \item[\textbf{б.}] Какой массив из элементов множества $\{1, 2, \dots, n\}$ содержит максимальное количество инверсий? Сколько инверсий в этом массиве?
    
    \item[\textbf{в.}] Какая существует взаимосвязь между временем сортировки методом вставок и количеством инверсий во входном массиве? Обоснуйте свой ответ.
    
    \item[\textbf{г.}] Разработайте алгоритм, определяющий количество инверсий, содержащихся в произвольной перестановке $n$ элементов, время работы которого в наихудшем случае равно $\Theta(n \lg n)$. (\textit{Указание:} модифицируйте алгоритм сортировки слиянием.)
\end{itemize}